\section{Unitary realizations}\label{sec:models}
\begin{table}[ht]\centering\small
\begin{tabular}{rrl}\toprule
$(\mu,j)$&Local index&Unitary realization\\\midrule
$(1,1)$&7&Tambara--Yamagami for $C_2^2$\\
$(1,3)$&9&$SU(2)_4$\\
$(1,4)$&14&$\Rep(D_{14})$\\
$(1,6)$&12&$\Rep(S_4)$\\
$(1,7)$&11&$PSU(2)_8$\\
$(1,10)$&15&$PSU(2)_9$\\
$(1,11)$&2&Tambara--Yamagami for $C_4$\\
$(1,13)$&4&The dimension-twelve model of \cite{LPR}\\
$(1,14)$&5&The dimension-twenty-four model of \cite{LPR}\\
$(1,16)$&1&$\Vec(C_5)$\\
$(2,27)$&33&$\Rep(A_5)$\\
$(2,36)$&5&$PSU(3)_4$\\
$(2,37)$&1&$\Rep(C_7\rtimes C_3)$\\
$(3,40)$&28&$C_2$-equivariantization of an Izumi $C_3+6$ category\\
$(3,75)$&8&$\Rep(\operatorname{AGL}(1,5))$\\
$(4,55)$&20&Izumi near-group $C_2^2+4$\\
$(4,123)$&6&Izumi near-group $C_4+4$\\
\bottomrule\end{tabular}
\caption{The seventeen positive rings. $D_{14}$ has order fourteen. Each tensor is constructed and matched by a unit-fixing permutation; the supplied permutation is checked.}\label{tab:models}\end{table}

The pointed, Tambara--Yamagami, finite-group and quantum-group entries
use their standard unitary realizations. The two nonstandard
multiplicity-one entries are the tensors of \cite[Section 3.5]{LPR},
with the unitary constructions in its Sections 4.2 and 4.4. The
$PSU(3)_4$ tensor agrees with the modular data in \cite{HR}.
The near-group realizations for $C_4+4$ and $C_2^2+4$ are
\cite[Examples 9.5--9.6]{Izumi}. These references establish existence,
not negative classifications used as an oracle.

The software constructs the $SU(2)$ tensors by the truncated
Clebsch--Gordan rule, the finite-group tensors by exact character
products, and the remaining displayed tensors in the cited conventions.
It checks their ring axioms, distinctness and a full tensor isomorphism
to the pinned census tensors. The unitary existence is supplied by the
constructions, not inferred just from passing the ring axioms.

For completeness, we explain the equivariantization entry, since an
order-two outer symmetry by itself would not supply the required action.
In \cite[Example 13.10]{Izumi}, the $C_3+6$ Cuntz realization has actual
symmetries generated by the pairs $(-1,R_{\pi/4})$ and
$(\id,L)$, where $L=\operatorname{diag}(1,-1)$ and $R_{\pi/4}$ is planar
rotation. Their product is $(-1,U)$ with $U=R_{\pi/4}L$ and $U^2=I$.
Thus it gives an actual order-two automorphism of the realization,
fixing the noninvertible endomorphism and the unit and interchanging
the two nontrivial invertibles. It fixes the unit inclusion in the square
of that endomorphism. On its six-dimensional self-fusion multiplicity
space $\CC[C_3]\otimes\CC^2$ it is inversion tensored with the
reflection $U$, and has trace zero. The two eigenspaces have dimension
three.

The $C_2$-equivariantization has simples $1,g,V,X,Y$:
two equivariant structures over the unit and the noninvertible object,
and one object over the orbit of the two other invertibles. The fixed
unit inclusion makes $X,Y$ self-dual. Tensoring equivariant structures
and using the two three-dimensional eigenspaces gives
\begin{align*}
 g^2&=1,&gV&=V,&gX&=Y,\\
 V^2&=1+g+V,&VX&=VY=X+Y,\\
 X^2=Y^2&=1+V+3X+3Y,&XY&=g+V+3X+3Y.
\end{align*}
The action is unitary, so its equivariantization is unitary. Its tensor
is generated from these products and matched exactly, just like the
other sixteen models.
